% GENERATED FILE. Edit canonical lessons, then run npm run generate:books.
% canonical-source-hash: fnv1a64:500d9c8f3e091aa8
% canonical-lessons: JA-C157-shokutaku, JA-C157-kabin, JA-C157-kazaru, JA-C157-kakeru, JA-C157-oku, JA-R157-review-152-157

\chapter{A Dining Table, a Vase and a Jacket on a Chair}
\label{ch:ja-a-dining-table-a-vase-and-a-jacket-on-a-chair}
\hlchaptermodality{157}

\begin{chapteropening}
\textbf{By the end of this chapter:} \emph{I can say where I put and hang things in a room.}

Complete the last lesson of chapter 157: I can say where I put and hang things in a room.
\end{chapteropening}

\section[\texorpdfstring{\ja{しょくたく} (shokutaku)}{しょくたく (shokutaku)}]{\ja{しょくたく} (shokutaku) — a dining table}

\label{lesson:JA-C157-shokutaku}



\begin{quote}
Before the new one: say the Japanese for to put in, then the Japanese for to arrange in a row.
\end{quote}

\subsection*{You'll want to know: \ja{しょくたく}}
\textbf{\ja{しょくたく}} — \emph{shokutaku} — \textquotedblleft{}a dining table\textquotedblright{}.

\begin{quote}
\textbf{\ja{かぞくと} \ja{しょくたくで} \ja{ごはんを} \ja{たべます。}} — \emph{kazoku to shokutaku de gohan o tabemasu} — \textquotedblleft{}I eat at the dining table with my family.\textquotedblright{}
\end{quote}

\begin{cousinweb}
\textbf{\ja{しょくたく}} (\emph{shokutaku}) is Sino-Japanese (kango): two Chinese characters read with their on'yomi, readings borrowed from Chinese. They are \emph{shoku}, \textquotedblleft{}eat, food\textquotedblright{} (Middle Chinese \emph{zyik}) and \emph{taku}, \textquotedblleft{}table, desk\textquotedblright{} (\emph{traewk}) (KANJIDIC2; Middle Chinese in Baxter's transcription).
\end{cousinweb}

\begin{grammarlens}[title={What you've built}]
One more word for everyday life at A2.
\end{grammarlens}

\subsection*{Your turn}
\begin{itemize}
\raggedright
  \item \emph{Say it:} \emph{shokutaku}
  \item \emph{Say it:} \emph{shokutaku}, once more
  \item \emph{Say it:} \emph{kazoku to shokutaku de gohan o tabemasu} — I eat at the dining table with my family.
  \item \emph{Recall:} say the Japanese for to put in, then the Japanese for to arrange in a row, then say \emph{shokutaku} again
\end{itemize}

\begin{tcolorbox}[breakable,colback=teal!4,colframe=teal!35!black,title={Before you move on}]
What does \ja{しょくたく} mean? (\textquotedblleft{}A dining table\textquotedblright{}.) Say it once more.
\end{tcolorbox}
\section[\texorpdfstring{\ja{かびん} (kabin)}{かびん (kabin)}]{\ja{かびん} (kabin) — a vase}

\label{lesson:JA-C157-kabin}



\begin{quote}
Before the new one: say the Japanese for to arrange in a row, then the Japanese for a dining table.
\end{quote}

\subsection*{You'll want to know: \ja{かびん}}
\textbf{\ja{かびん}} — \emph{kabin} — \textquotedblleft{}a vase\textquotedblright{}.

\begin{quote}
\textbf{\ja{しろい} \ja{かびんは} \ja{たかいです。}} — \emph{shiroi kabin wa takai desu} — \textquotedblleft{}The white vase is expensive.\textquotedblright{}
\end{quote}

\begin{cousinweb}
\textbf{\ja{かびん}} (\emph{kabin}) is Sino-Japanese (kango): two Chinese characters read with their on'yomi, readings borrowed from Chinese. They are \emph{ka}, \textquotedblleft{}flower\textquotedblright{} (Middle Chinese \emph{xwae}) and \emph{bin}, \textquotedblleft{}bottle, jar\textquotedblright{} (\emph{beng}) (KANJIDIC2; Middle Chinese in Baxter's transcription).
\end{cousinweb}

\begin{grammarlens}[title={What you've built}]
One more word for everyday life at A2.
\end{grammarlens}

\subsection*{Your turn}
\begin{itemize}
\raggedright
  \item \emph{Say it:} \emph{kabin}
  \item \emph{Say it:} \emph{kabin}, once more
  \item \emph{Say it:} \emph{shiroi kabin wa takai desu} — The white vase is expensive.
  \item \emph{Recall:} say the Japanese for to arrange in a row, then the Japanese for a dining table, then say \emph{kabin} again
\end{itemize}

\begin{tcolorbox}[breakable,colback=teal!4,colframe=teal!35!black,title={Before you move on}]
What does \ja{かびん} mean? (\textquotedblleft{}A vase\textquotedblright{}.) Say it once more.
\end{tcolorbox}
\section[\texorpdfstring{\ja{かざる} (kazaru)}{かざる (kazaru)}]{\ja{かざる} (kazaru) — to decorate; to put on display}

\label{lesson:JA-C157-kazaru}



\begin{quote}
Before the new one: say the Japanese for a dining table, then the Japanese for a vase.
\end{quote}

\subsection*{You'll want to know: \ja{かざる}}
\textbf{\ja{かざる}} — \emph{kazaru} — \textquotedblleft{}to decorate; to put on display\textquotedblright{}.

\begin{quote}
\textbf{\ja{げんかんに} \ja{はなを} \ja{かざります。}} — \emph{genkan ni hana o kazarimasu} — \textquotedblleft{}I display flowers in the entry hall.\textquotedblright{}
\end{quote}

Polite \textbf{\ja{かざります}}, \emph{kazarimasu}; a godan verb.

\begin{cousinweb}
\textbf{\ja{かざる}} (\emph{kazaru}) is a native Japanese word (wago), inherited rather than borrowed from Chinese; its kanji was attached later as a kun'yomi, and its older origin is not settled (JMdict, KANJIDIC2).
\end{cousinweb}

\begin{grammarlens}[title={What you've built}]
One more word for everyday life at A2.
\end{grammarlens}

\subsection*{Your turn}
\begin{itemize}
\raggedright
  \item \emph{Say it:} \emph{kazaru}
  \item \emph{Say it:} \emph{kazaru}, once more
  \item \emph{Say it:} \emph{genkan ni hana o kazarimasu} — I display flowers in the entry hall.
  \item \emph{Recall:} say the Japanese for a dining table, then the Japanese for a vase, then say \emph{kazaru} again
\end{itemize}

\begin{tcolorbox}[breakable,colback=teal!4,colframe=teal!35!black,title={Before you move on}]
What does \ja{かざる} mean? (\textquotedblleft{}To decorate; to put on display\textquotedblright{}.) Say it once more.
\end{tcolorbox}
\section[\texorpdfstring{\ja{かける} (kakeru)}{かける (kakeru)}]{\ja{かける} (kakeru) — to hang; to hang up}

\label{lesson:JA-C157-kakeru}



\begin{quote}
Before the new one: say the Japanese for a vase, then the Japanese for to decorate.
\end{quote}

\subsection*{You'll want to know: \ja{かける}}
\textbf{\ja{かける}} — \emph{kakeru} — \textquotedblleft{}to hang; to hang up\textquotedblright{}.

\begin{quote}
\textbf{\ja{うわぎを} \ja{いすに} \ja{かけます。}} — \emph{uwagi o isu ni kakemasu} — \textquotedblleft{}I hang my jacket on the chair.\textquotedblright{}
\end{quote}

Polite \textbf{\ja{かけます}}, \emph{kakemasu}; an ichidan verb.

\begin{cousinweb}
\textbf{\ja{かける}} (\emph{kakeru}) is a native Japanese word (wago), inherited rather than borrowed from Chinese; its kanji was attached later as a kun'yomi, and its older origin is not settled (JMdict, KANJIDIC2).
\end{cousinweb}

\begin{grammarlens}[title={What you've built}]
One more word for everyday life at A2.
\end{grammarlens}

\subsection*{Your turn}
\begin{itemize}
\raggedright
  \item \emph{Say it:} \emph{kakeru}
  \item \emph{Say it:} \emph{kakeru}, once more
  \item \emph{Say it:} \emph{uwagi o isu ni kakemasu} — I hang my jacket on the chair.
  \item \emph{Recall:} say the Japanese for a vase, then the Japanese for to decorate, then say \emph{kakeru} again
\end{itemize}

\begin{tcolorbox}[breakable,colback=teal!4,colframe=teal!35!black,title={Before you move on}]
What does \ja{かける} mean? (\textquotedblleft{}To hang; to hang up\textquotedblright{}.) Say it once more.
\end{tcolorbox}
\section[\texorpdfstring{\ja{おく} (oku)}{おく (oku)}]{\ja{おく} (oku) — to put; to place}

\label{lesson:JA-C157-oku}



\begin{quote}
Before the new one: say the Japanese for to decorate, then the Japanese for to hang.
\end{quote}

\subsection*{You'll want to know: \ja{おく}}
\textbf{\ja{おく}} — \emph{oku} — \textquotedblleft{}to put; to place\textquotedblright{}.

\begin{quote}
\textbf{\ja{かびんを} \ja{しょくたくに} \ja{おきます。}} — \emph{kabin o shokutaku ni okimasu} — \textquotedblleft{}I put the vase on the dining table.\textquotedblright{}
\end{quote}

Polite \textbf{\ja{おきます}}, \emph{okimasu}; a godan verb. Spelled and said like \textbf{\ja{おきます}}, \emph{okimasu}, \textquotedblleft{}I get up\textquotedblright{}; here it takes an object with \textbf{\ja{を}}.

\begin{cousinweb}
\textbf{\ja{おく}} (\emph{oku}) is a native Japanese word (wago), inherited rather than borrowed from Chinese; its kanji was attached later as a kun'yomi, and its older origin is not settled (JMdict, KANJIDIC2).
\end{cousinweb}

\begin{grammarlens}[title={What you've built}]
One more word for everyday life at A2.
\end{grammarlens}

\subsection*{Your turn}
\begin{itemize}
\raggedright
  \item \emph{Say it:} \emph{oku}
  \item \emph{Say it:} \emph{oku}, once more
  \item \emph{Say it:} \emph{kabin o shokutaku ni okimasu} — I put the vase on the dining table.
  \item \emph{Recall:} say the Japanese for to decorate, then the Japanese for to hang, then say \emph{oku} again
\end{itemize}

\begin{tcolorbox}[breakable,colback=teal!4,colframe=teal!35!black,title={Before you move on}]
What does \ja{おく} mean? (\textquotedblleft{}To put; to place\textquotedblright{}.) Say it once more.
\end{tcolorbox}
\section[(dialogue)]{Words from chapters 152–157 — a review}

\label{lesson:JA-R157-review-152-157}



\begin{quote}
Say the Japanese for preparations and to hurry.
\end{quote}

\subsection*{Your turn}
Say these aloud:

\begin{itemize}
\raggedright
  \item preparations, to hurry, to go out, to forget, to be late
  \item going out, coming home, something left behind, to lose, to be in time
  \item usually, a habit, everyday life, to get by, to get used to
  \item to move house, to carry, to shift, furniture, the front entrance
  \item a closet, a drawer, a home office, to put in, to arrange in a row
  \item a dining table, a vase, to decorate, to hang, to put
\end{itemize}

\begin{tcolorbox}[breakable,colback=teal!4,colframe=teal!35!black,title={Before you move on}]
Preparations? (\textbf{\ja{したく}}, \emph{shitaku}.) To hurry? (\textbf{\ja{いそぐ}}, \emph{isogu}.) To go out? (\textbf{\ja{でかける}}, \emph{dekakeru}.) To forget? (\textbf{\ja{わすれる}}, \emph{wasureru}.) To be late? (\textbf{\ja{おくれる}}, \emph{okureru}.) Going out? (\textbf{\ja{がいしゅつ}}, \emph{gaishutsu}.) Coming home? (\textbf{\ja{きたく}}, \emph{kitaku}.) Something left behind? (\textbf{\ja{わすれもの}}, \emph{wasuremono}.) To lose? (\textbf{\ja{なくす}}, \emph{nakusu}.) To be in time? (\textbf{\ja{まにあう}}, \emph{maniau}.) Usually? (\textbf{\ja{ふだん}}, \emph{fudan}.) A habit? (\textbf{\ja{くせ}}, \emph{kuse}.) Everyday life? (\textbf{\ja{にちじょう}}, \emph{nichijō}.) To get by? (\textbf{\ja{くらす}}, \emph{kurasu}.) To get used to? (\textbf{\ja{なれる}}, \emph{nareru}.) To move house? (\textbf{\ja{ひっこす}}, \emph{hikkosu}.) To carry? (\textbf{\ja{はこぶ}}, \emph{hakobu}.) To shift? (\textbf{\ja{うごかす}}, \emph{ugokasu}.) Furniture? (\textbf{\ja{かぐ}}, \emph{kagu}.) The front entrance? (\textbf{\ja{げんかん}}, \emph{genkan}.) A closet? (\textbf{\ja{おしいれ}}, \emph{oshiire}.) A drawer? (\textbf{\ja{ひきだし}}, \emph{hikidashi}.) A home office? (\textbf{\ja{しょさい}}, \emph{shosai}.) To put in? (\textbf{\ja{いれる}}, \emph{ireru}.) To arrange in a row? (\textbf{\ja{ならべる}}, \emph{naraberu}.) A dining table? (\textbf{\ja{しょくたく}}, \emph{shokutaku}.) A vase? (\textbf{\ja{かびん}}, \emph{kabin}.) To decorate? (\textbf{\ja{かざる}}, \emph{kazaru}.) To hang? (\textbf{\ja{かける}}, \emph{kakeru}.) To put? (\textbf{\ja{おく}}, \emph{oku}.)
\end{tcolorbox}
